\documentclass[12pt,a4paper]{report}
\usepackage[utf8]{inputenc}
\usepackage[T1]{fontenc}
\usepackage{graphicx}
\usepackage{geometry}
\usepackage[output=latex]{plantuml}
% PlantUML's TikZ boxes assume 10pt text; do not inherit the report's 12pt.
\renewcommand{\PlantUmlTikzFont}{\sffamily\fontsize{10}{12}\selectfont}
\geometry{hmargin=2.5cm,vmargin=2.5cm}
\usepackage[hidelinks]{hyperref}
\usepackage[figure]{hypcap}
\usepackage{titlesec}
\usepackage{tabularx}
\usepackage{amsmath}
\usepackage{enumitem}
\usepackage{float}

\begin{document}

\title{Report \newline Software Design}
\author{Nguyen Hoang Nhan LE \and Quynh Anh HOANG \and Hanane }

% --- CUSTOM UT CAPITOLE COVER PAGE ---
\hypersetup{pageanchor=false}
\begin{titlepage}
    \centering
    \begin{minipage}{0.4\textwidth}
        % Make sure to download the official UT Capitole logo from the intranet
        %\includegraphics[width=\textwidth]{logo_ut_capitole.png}
        \textbf{\begin{center}
            UNIVERSITÉ TOULOUSE CAPITOLE
        \end{center}}
    \end{minipage}
    \hfill
    \begin{minipage}{0.4\textwidth}
        \raggedleft
        \textit{Année Universitaire 2026-2027}
    \end{minipage}
    
    \vfill
    {\large \textbf{REPORT}} \\[1cm]
    {\huge \bfseries Software Design } \\[0.5cm]
    {\large Assignment 1}
    
    \vfill
    \begin{minipage}{0.4\textwidth}
        \textbf{Group member :} \\
        Nguyen Hoang Nhan \textsc{LE} \\
        Hanane \textsc{ELOUCHEFOUNE}
    \end{minipage}
    \hfill
    \begin{minipage}{0.4\textwidth}
        \raggedleft
        \textbf{Professor :} \\
        Prof. Laurent \textsc{Perrussel}
    \end{minipage}
    
    \vfill
    {\large Date de soutenance : 30/10/2026}
\end{titlepage}
\hypersetup{pageanchor=true}

% --- MANDATORY UNIVERSITY DISCLAIMER ---
% \thispagestyle{empty}
% \noindent\textit{« L'Université n'entend donner aucune approbation ni approbation aux opinions émises dans les thèses ou rapports ; ces opinions doivent être considérées comme propres à leurs auteurs. »}
% \clearpage

% --- PRELIMINARY PAGES ---
\tableofcontents
\listoffigures
\listoftables
\clearpage

% --- MAIN CONTENT ---
\chapter{Usecase Digrams}
\section{Indicate usecases and actors}
From the description of the application, we defined the main actors
of our application.
\begin{table}[h]
    \centering
    \begin{tabularx}{\textwidth}{|l|X|}
        \hline
        \textbf{Actors} & \textbf{Reason} \\
        \hline
        Customer & As mentioned in the assignment: ``...customer starts a reservation by mentioning where they want to pick up the car...'' \\
        % \hline
        % Payment Service Provider & Hanlding payment by credit card is a
        % quite sensitive operation, website usually offload these task to a third party provider specialized in
        % payment processing. \\
        \hline
        Agency Employee & As mentioned in the assingment: ``Based on the
    chosen category, the agency employee will allocate a specific car to the customer.'' \\
        \hline
        Hotline Operator & As metioned in the assignment: ``The 24/7 hotline will contact her back and organize with her the repairing actions'' \\
        \hline
    \end{tabularx}
    \caption{Actors in system}
\end{table}

\textbf{Design choices}: Each represent a role interacting with the system.
We choosed to design the system with one Customer
for the actors before and after signing in because the action of
authentication change the customer's access to the system, not changing
their role. Agency Employee represent the employee working at the agency
location and handle pickup and return action. Hotline operator handle
report and coordinate assistances following an incident.

These actors with participate in several usecase:

\begin{center}
\begin{table}[h]
    \begin{tabularx}{\textwidth}{|l|X|}
        \hline
        \textbf{Usecase} & \textbf{Participation} \\
        \hline
        Make reservation & Customer \\
        \hline
        Pay & Customer\\
        \hline
        Cancel reservation & Customer \\
        \hline
        Pickup car & Customer, Agency Employee \\
        \hline
        Return car & Customer, Agency Employee \\
        \hline
        Manage incident & Customer, Hotline Operator \\
        \hline
        Report incident & Customer \\
        \hline
    \end{tabularx}
\caption{List of usecases and their participations}
\end{table}
\end{center}
% \section{Usecases details}
% \subsection{Make reservation}

% \begin{table}[H]
%     \centering
%     \begin{tabularx}{\textwidth}{|l|X|}
%         \hline
%         Name: & Make Reservation \\
%         \hline
%         Short description: & A customer reserve a car from the website of the company. \\
%         \hline
%         Precondition: & None \\
%         \hline
%         Postcondition: & A car booking is made; payment is completed for prepay. \\
%         \hline
%         Error situations: & None\\
%         \hline
%         System state in the event of an error: & No confirmed booking is created\\
%         \hline
%         Actors: & Customer \\
%         \hline
%         Trigger: & Customer wanting to rent a car. \\
%         \hline
%         Standard process: & \begin{enumerate}[nosep, leftmargin=*,label=(\arabic*)]
%             \item Customer goes to the website and starts a reservation.
%             \item Customer enters the pickup location and pickup and return date.
%             \item The system presents car category and their prepay and pay-later prices.
%             \item Customer select the car category.
%             \item Customer select the payment option.
%             \item Customer optionally choose soome extra insurances and options.
%             \item Customer signs in to continue the booking process.
%             \item Customer provides provides their email, first and last names, and address: street, city, country, and zip code.
%             \item If the customer has chose the prepay option, they then pay with credit card.
%             \item The system records and confirms the booking.
%         \end{enumerate} \\
%         \hline
%     \end{tabularx}
%     \caption{Make reservation usecase details}
% \end{table}

% \subsection{Pay}
% \begin{table}[H]
%     \centering
%     \begin{tabularx}{\textwidth}{|l|X|}
%         \hline
%         Name: & Pay \\
%         \hline
%         Short description: & A customer do a transaction to transfer an amount of money to the company \\
%         \hline
%         Precondition: & None \\
%         \hline
%         Postcondition: & The rental company recieves the money \\
%         \hline
%         Error situations: & None\\
%         \hline
%         System state in the event of an error: & No amount of money is subtract from the customer account.
%         The company does not receive the money. \\
%         \hline
%         Actors: & Customer \\
%         \hline
%         Trigger: & Customer wanting to rent a car or canceling a prepaid reservation. \\
%         \hline
%         Standard process: & \begin{enumerate}[nosep, leftmargin=*,label=(\arabic*)]
%             \item Customer goes to the website and starts a reservation.
%             \item Customer enters the pickup location and pickup and return date.
%             \item The system presents car category and their prepay and pay-later prices.
%             \item Customer select the car category.
%             \item Customer select the payment option.
%             \item Customer optionally choose soome extra insurances and options.
%             \item Customer signs in to continue the booking process.
%             \item Customer provides provides their email, first and last names, and address: street, city, country, and zip code.
%             \item If the customer has chose the prepay option, they then pay with credit card.
%             \item The system records and confirms the booking.
%         \end{enumerate} \\
%         \hline
%     \end{tabularx}
%     \caption{Make reservation usecase details}
% \end{table}
% \subsection{Cancel reservation}
% \subsection{Pickup car}
% \subsection{Manage incident}
% \subsection{Report incident}

\clearpage
\section{Usecase diagram}

\begin{center}
% Reserve space for the section heading and keep the caption outside scaling.
\begin{adjustbox}{max totalsize={\linewidth}{0.8\textheight}}
\plantumlinput{digrams/usecase/usecase.pu}
\end{adjustbox}
\captionof{figure}{General usecase diagram}
\label{fig:usecase-diagram}
\end{center}

\section{Detail of use cases associated to the action "pick up car"}
\subsection{System sequence diagram}

\begin{figure}[H]
    \centering
    \begin{adjustbox}{
        max totalsize={\textwidth}{0.75\textheight},
        center
    }
    \plantumlinput{digrams/sequence/pickup-with-condition-ssd.pu}
    \end{adjustbox}
    \caption{System sequence diagram for pickup car action}
\end{figure}

\subsection{Lo-Fi interface}
We have designed a lo-fi interface serving the pickup action. The activities follow the sequence:
find reservation,
review it details,
allocate a car from the booked category,
check and scan the driver license,
record inspection result with customer's digital signature.

You can click on the navigation button link the the corresponding figures in this PDF that help 
illustrate the workflow steps. Every actions must succeed before going to the next step and the
key is only handed out after all checks and payment are complete.

\begin{figure}[H]
    \centering
    \begin{adjustbox}{
        max totalsize={\textwidth}{0.75\textheight},
        center
    }
    \plantumlinput{"digrams/interface/pickup-main.pu"}
    \end{adjustbox}
    \caption{Main screen of Lo-fi inteface for pickup action}
    \label{fig:lofi-ui-main}
\end{figure}
\begin{figure}[H]
    \centering
    \begin{adjustbox}{
        max totalsize={\textwidth}{0.75\textheight},
        center
    }
    \plantumlinput{"digrams/interface/pickup-open-reservation.pu"}
    \end{adjustbox}
    \caption{Reservation details for the pickup action}
    \label{fig:lofi-ui-reservation-details}
\end{figure}

\begin{figure}[H]
    \centering
    \begin{adjustbox}{max totalsize={\textwidth}{0.75\textheight},center}
    \plantumlinput{"digrams/interface/pickup-car-allocation.pu"}
    \end{adjustbox}
    \caption{Select and allocate an available car in the booked category}
    \label{fig:lofi-ui-car-allocation}
\end{figure}

\begin{figure}[H]
    \centering
    \begin{adjustbox}{max totalsize={\textwidth}{0.75\textheight},center}
    \plantumlinput{"digrams/interface/pickup-driving-licence.pu"}
    \end{adjustbox}
    \caption{Check the driving licence and store its scan}
    \label{fig:lofi-ui-driving-licence}
\end{figure}

\begin{figure}[H]
    \centering
    \begin{adjustbox}{max totalsize={\textwidth}{0.75\textheight},center}
    \plantumlinput{"digrams/interface/pickup-inspection.pu"}
    \end{adjustbox}
    \caption{Record the joint inspection and the customer's digital signature}
    \label{fig:lofi-ui-inspection}
\end{figure}

\begin{figure}[H]
    \centering
    \begin{adjustbox}{max totalsize={\textwidth}{0.75\textheight},center}
    \plantumlinput{"digrams/interface/pickup-payment.pu"}
    \end{adjustbox}
    \caption{Record payment for a pay-at-pickup reservation}
    \label{fig:lofi-ui-payment}
\end{figure}

\begin{figure}[H]
    \centering
    \begin{adjustbox}{max totalsize={\textwidth}{0.75\textheight},center}
    \plantumlinput{"digrams/interface/pickup-handover.pu"}
    \end{adjustbox}
    \caption{Review completed checks before handing over the keys}
    \label{fig:lofi-ui-handover}
\end{figure}

\chapter{Class Diagram}

Class digram represents the main classes that is related to the primary business of the application.
From the assignment text, we design the system with the following classes:
\begin{itemize}
    \item \textbf{Customer} class represent the 
\end{itemize}
\begin{figure}[H]
    \centering
    % Fit both dimensions, leaving room for the chapter heading and caption.
    % \begin{adjustbox}{max totalsize={\linewidth}{0.72\textheight}}
    \plantumlinput{digrams/class/class.pu}
    % \end{adjustbox}
    \captionof{figure}{Class diagram}
    \label{fig:class-diagram}
\end{figure}


\chapter{Sequence Diagram}
\section{Reservation action}
\begin{figure}[H]
    \centering
    \plantumlinput{"digrams/sequence/reservation.pu"}
    \caption{Sequence diagram for reservation}
\end{figure}
\end{document}
